\documentclass[12pt,a4paper]{article}

\usepackage[margin=1in]{geometry}
\usepackage[shortlabels]{enumitem}
\usepackage{hyperref}

\hypersetup{
  colorlinks=true,
  linkcolor=blue,
  urlcolor=blue,
}

\setlength{\parskip}{0.6em}
\setlength{\parindent}{0pt}

\title{Speaker Notes:\\
  ML Surrogate for SF\textsubscript{6} Plasma Discharge Modeling}
\author{Progress Update Script --- Internal Use}
\date{April 2026}

\begin{document}
\maketitle
\tableofcontents
\clearpage

%% ============================================================
\section*{How to Use This Document}

These notes are written as a progress-update script for a supervisor
who already knows the 2D TEL finite-difference solver from the
previous meeting. Everything here is \emph{new work since that
meeting}. The speaking tone is informal and first-person: ``Last time
we discussed X. Since then, I have done Y.''

Each slide section contains:
\begin{itemize}
  \item \textbf{Purpose} --- what the slide must accomplish.
  \item \textbf{What to Say} --- 2--4 paragraphs of natural first-person speech.
  \item \textbf{Figure Guidance} --- how to read the plot or diagram on screen.
  \item \textbf{Likely Questions} --- what the supervisor may ask and how to handle it.
  \item \textbf{Transition} --- the bridge sentence to the next slide.
\end{itemize}

Speak at roughly 2 minutes per slide; the full update should land at
35--40 minutes plus Q\&A.

\clearpage

%% ============================================================
\section{Slide 1 --- Title}
\label{sec:slide1}

\subsection*{Purpose}
Frame this as a progress update, not a fresh introduction.
The supervisor already knows the problem and the solver; establish
that today's content is new work.

\subsection*{What to Say}

Thanks for making time. Last time we met, I walked through the 2D TEL
finite-difference solver --- the governing equations, the mesh, and why
the runtime was a bottleneck. Today I want to show you everything
that has happened since then.

The headline result is an ensemble surrogate that hits a
log\textsubscript{10} fluorine-density RMSE of 0.00154 --- roughly
0.35\% error in linear density --- with speedups of 531$\times$ to
874$\times$ over the solver. But the path to get there involved some
dead ends that I think are just as worth discussing. I will cover
those honestly.

\subsection*{Figure Guidance}
No figure on this slide. Title and date should be visible.

\subsection*{Likely Questions}
None expected here. If asked ``Is this a PINN talk?'' --- answer:
PINNs were one of the first things I tried; I will show why they
failed for this problem.

\subsection*{Transition}
Let me start with a quick recap of where we left off.

\clearpage

%% ============================================================
\section{Slide 2 --- Recap: Where We Left Off}
\label{sec:slide2}

\subsection*{Purpose}
One slide, thirty seconds. Remind the supervisor of the solver,
the inputs, and the output we care about. Do not re-derive
anything --- just anchor the conversation.

\subsection*{What to Say}

Quick recap so we are on the same page. We have the 2D TEL
drift-diffusion solver for SF\textsubscript{6} discharges. It takes
operating conditions --- pressure, voltage, gap distance --- and
produces spatial profiles of species densities on a non-uniform mesh.
The quantity we care most about is the steady-state atomic fluorine
density, nF, because it drives component degradation in switchgear.
The solver is validated but slow, which motivated the surrogate work
I am about to describe.

You already know all of this. I just want to make sure we are
starting from the same place before I show the new directions.

\subsection*{Figure Guidance}
If a schematic of the simulation domain is shown, gesture to it
briefly. Do not explain it --- the supervisor has seen it before.

\subsection*{Likely Questions}
\begin{itemize}
  \item \emph{Has anything changed in the solver since last time?}
    --- No changes to the solver itself; all the new work is on the
    ML side and the LXCat cross-section integration.
\end{itemize}

\subsection*{Transition}
Since that meeting, I have been exploring three directions. Let me
give you the overview before diving into each.

\clearpage

%% ============================================================
\section{Slide 3 --- Three Directions Overview}
\label{sec:slide3}

\subsection*{Purpose}
Give the supervisor a roadmap for the rest of the talk. Three
directions: (1) PINNs, (2) supervised surrogates, (3) LXCat
cross-section integration. Set expectations for the narrative arc.

\subsection*{What to Say}

Since our last meeting, I have pursued three main threads. First, I
tried physics-informed neural networks --- PINNs --- because they are
the obvious first idea when you have a PDE solver. That did not work,
and I will explain why. Second, I pivoted to a purely supervised
surrogate approach, which is where all the positive results came from.
Third, I integrated updated cross-section data from LXCat into the
solver and studied how it affects both the physics and the surrogate.

The story arc is: PINN attempt, failure, pivot, then a systematic
progression through four surrogate versions that brought the error
down by two orders of magnitude. Along the way, the LXCat work
produced a genuine surprise about dataset sensitivity. I will also
cover architecture sweeps, ablation studies, and several negative
results that I think are important to document.

\subsection*{Figure Guidance}
If a roadmap or flowchart is shown, point to the three branches.
Emphasise that the PINN branch is a dead end --- the supervisor
should not expect positive results from it.

\subsection*{Likely Questions}
\begin{itemize}
  \item \emph{Why did you start with PINNs?}
    --- They are the most-cited approach in the literature for
    PDE surrogates, so it was the natural starting point.
  \item \emph{How long did the PINN work take before you pivoted?}
    --- Be honest about the time spent.
\end{itemize}

\subsection*{Transition}
Let me start with the PINN attempt.

\clearpage

%% ============================================================
\section{Slide 4 --- PINN Attempt and Failure}
\label{sec:slide4}

\subsection*{Purpose}
Report the PINN experiments honestly. Two variants tried, both
failed. This is a valuable negative result that justifies the
pivot to supervised learning.

\subsection*{What to Say}

PINNs were the first thing I tried after our last meeting.
I tested two variants. The data-only PINN --- essentially a neural
network trained on solver data but using the PINN framework with
automatic differentiation and collocation points --- does
eventually converge, but it is 23,000$\times$ slower than standard
supervised training to reach the same loss. The framework overhead
is enormous and provides no benefit when you already have labelled
data from the solver.

The second variant was a PDE-constrained PINN, where I added the
governing equations as a soft penalty in the loss. This one
diverges outright. The PDE residual and the data residual compete,
and the optimiser cannot satisfy both simultaneously. I tried
adaptive loss weighting, curriculum strategies, and NTK-based
balancing --- none of them stabilised training.

The takeaway is straightforward: PINNs are not the right tool
here. We have abundant, trusted solver data, so there is no reason
to pay the overhead of physics-informed training. This failure is
what motivated the pivot to a purely supervised approach, which I
will cover next.

\subsection*{Figure Guidance}
If loss curves are shown, point to the data-only PINN curve (slowly
descending) versus the standard training curve (converging orders of
magnitude faster). Point to the PDE-constrained curve diverging or
oscillating. The visual contrast makes the argument immediately.

\subsection*{Likely Questions}
\begin{itemize}
  \item \emph{Is the PDE too stiff for PINNs?}
    --- Likely yes; the SF\textsubscript{6} chemistry has rate
    constants spanning many orders of magnitude.
  \item \emph{Could a better PINN implementation work?}
    --- Possibly, but the 23,000$\times$ overhead for the data-only
    case suggests the framework itself is the bottleneck. Even if
    PDE-constrained training could be stabilised, we would still
    face that overhead.
  \item \emph{Did you try DeepONet or Fourier Neural Operator?}
    --- Not yet; those are possible future directions, but the
    supervised approach already works well.
\end{itemize}

\subsection*{Transition}
So PINNs were a dead end. That led me to ask: what if we just do
straightforward supervised learning?

\clearpage

%% ============================================================
\section{Slide 5 --- Pivot to Supervised Surrogates}
\label{sec:slide5}

\subsection*{Purpose}
Explain the supervised surrogate strategy: generate training data
from the solver, train a neural network to map inputs to spatial
profiles. Justify working in log space.

\subsection*{What to Say}

After the PINN failure, I went back to basics. The idea is simple:
run the solver across a grid of input conditions, collect the nF
profiles, and train a feedforward network to reproduce them. No
physics in the loss --- just mean-squared error on
log\textsubscript{10} nF.

Working in log space turned out to be critical. The fluorine density
spans several orders of magnitude across the spatial domain, so
predicting in linear space forces the network to balance accuracy at
high and low densities, which it does poorly. In log space, the
dynamic range is compressed and the regression problem is
well-conditioned. This single transformation is probably the most
important design choice in the entire pipeline.

The output is not a single number --- it is the full spatial profile
at every mesh node. So the network has a multi-dimensional output,
and the architecture needs to handle that gracefully. That is where
the surrogate versioning comes in: I went through four major versions,
each addressing a specific shortcoming of the previous one.

\subsection*{Figure Guidance}
If a schematic of the input--output mapping is shown, point to the
inputs on the left and the spatial profile on the right. Emphasise
that this is a curve prediction problem, not scalar regression.

\subsection*{Likely Questions}
\begin{itemize}
  \item \emph{How many solver runs are in the training set?}
    --- Give the exact number; mention 80/10/10 train/val/test split.
  \item \emph{Why not predict in linear space with a log loss?}
    --- Equivalent in principle, but log-space MSE is simpler and
    gave better results empirically.
\end{itemize}

\subsection*{Transition}
Let me walk you through the four versions.

\clearpage

%% ============================================================
\section{Slide 6 --- Surrogate v1 through v4 Progression}
\label{sec:slide6}

\subsection*{Purpose}
Show the error progression across four model versions, each
incorporating a specific improvement. This demonstrates systematic
iteration, not luck.

\subsection*{What to Say}

Here is the progression. Version 1 was a baseline MLP --- nothing
fancy, just a few hidden layers with ReLU activations. It achieved
an RMSE of 0.133, which is terrible --- more than 30\% error in
linear density. It was a sanity check that the pipeline works, but
the model itself was not useful.

Version 2 added better input normalisation and a wider network,
bringing the error down to 0.111. Still bad, but a 17\% relative
improvement. Version 3 is where things got interesting: I introduced
bias initialisation to the training-set mean and switched to a
cosine-annealing learning rate schedule. That dropped the RMSE to
0.023 --- nearly a 5$\times$ improvement over v2. This confirmed
that training recipe matters more than raw architecture at this stage.

Version 4 introduced the separate-heads architecture and weight decay.
The RMSE fell to 0.003. So across four iterations, we went from 0.133
to 0.003 --- roughly a 44$\times$ improvement. Each version targeted
a specific bottleneck, and the progression was monotonic. No lucky
accidents.

\subsection*{Figure Guidance}
If a bar chart or waterfall showing v1 through v4 errors is
displayed, walk through each bar left to right. Emphasise the big
jump from v2 to v3 (bias init + LR schedule) and from v3 to v4
(separate heads).

\subsection*{Likely Questions}
\begin{itemize}
  \item \emph{What changed between each version?}
    --- Summarise: v1$\to$v2 was normalisation/width, v2$\to$v3 was
    bias init + LR schedule, v3$\to$v4 was separate heads + weight
    decay.
  \item \emph{Why not skip to v4?}
    --- Each version taught us something; the ablation study (later)
    confirms which ingredients matter.
\end{itemize}

\subsection*{Transition}
Version 4 got us to 0.003. But can we do better? Let me show the
accuracy analysis.

\clearpage

%% ============================================================
\section{Slide 7 --- Surrogate v4 Accuracy}
\label{sec:slide7}

\subsection*{Purpose}
Deep-dive on v4 performance. Show where it succeeds and where
residual errors remain, motivating the ensemble and the architecture
sweep.

\subsection*{What to Say}

Version 4 achieves an RMSE of 0.00292 on the test set. In linear
terms, that is about 0.67\% error --- already quite good for a
surrogate. But when I looked at the spatial error distribution, the
errors are not uniform. The near-electrode regions, where the nF
gradient is steepest, have higher residuals than the bulk.

This told me two things. First, a single model with shared output
weights is struggling to simultaneously capture the steep-gradient
and flat-gradient regions. That is what motivated the separate-heads
design in v4, and it helped, but there was still room for improvement.
Second, seed-to-seed variance was noticeable --- different random
initialisations gave different error patterns. That pointed toward
ensembling as the next step.

So the path forward was clear: run an architecture sweep to find the
best head configuration, then ensemble the winners. That is exactly
what I did.

\subsection*{Figure Guidance}
If a spatial error heatmap or profile overlay is shown, point to the
electrode regions where the residual is highest. If a parity plot is
shown, point to the slight scatter at the distribution tails.

\subsection*{Likely Questions}
\begin{itemize}
  \item \emph{Is 0.00292 already good enough for engineering use?}
    --- Arguably yes, but we can do better, and the ensemble also
    gives us uncertainty quantification.
  \item \emph{How much does seed variance matter in practice?}
    --- Enough to change the third significant figure of RMSE, which
    is why ensembling is worth the cost.
\end{itemize}

\subsection*{Transition}
Before the ensemble, I integrated the LXCat cross-section data.
Let me show you what that revealed.

\clearpage

%% ============================================================
\section{Slide 8 --- LXCat Integration}
\label{sec:slide8}

\subsection*{Purpose}
Explain why LXCat data was integrated and what it means for the
surrogate strategy. This is new work since the last meeting.

\subsection*{What to Say}

One thing I wanted to investigate since our last discussion was how
sensitive the results are to the choice of electron-scattering
cross-sections. Our solver has always used a legacy hand-curated
dataset. LXCat is the community-maintained open-access database that
the plasma physics community increasingly treats as the reference
source. So I swapped in the LXCat cross-sections and re-ran the full
simulation matrix.

Integrating LXCat on the solver side was straightforward --- I
replaced the lookup tables. But it doubled the simulation effort
because every input condition now had to be run with both datasets.
The payoff is that we get paired data: for each input condition, we
can compare legacy vs.\ LXCat outputs directly and ask whether the
surrogate needs to distinguish between them.

The answer turned out to be surprising, and I will get to that
in a moment. First, let me show you what actually changed
in the physics.

\subsection*{Figure Guidance}
If a table of cross-section sources is shown, point to the specific
datasets used. If a workflow diagram shows the solver with swappable
cross-section tables, highlight the swap point.

\subsection*{Likely Questions}
\begin{itemize}
  \item \emph{Which LXCat dataset specifically?}
    --- Name the source and version.
  \item \emph{Did you re-validate the solver with LXCat data?}
    --- Yes, coming on the solver validation slide.
\end{itemize}

\subsection*{Transition}
Here is what changed in the physics.

\clearpage

%% ============================================================
\section{Slide 9 --- LXCat Physics Impact}
\label{sec:slide9}

\subsection*{Purpose}
Show that LXCat cross-sections produce real, measurable changes in
intermediate quantities (Te, ne), even though the effect on nF is
small.

\subsection*{What to Say}

When I compared the legacy and LXCat solver outputs, the intermediate
quantities shifted substantially. Electron temperature increased by
56\% under LXCat cross-sections. Electron density decreased by 69\%.
These are large, physically meaningful changes --- the two datasets
really do produce different plasma states.

But here is the surprising part: fluorine density --- the quantity we
actually predict --- changed by less than 2\%. The big shifts in Te
and ne largely cancel out when propagated through the full chemistry
to nF. Fluorine density is an integrated downstream quantity that
depends on many competing reactions, and errors in individual rates
partially offset each other.

This was not obvious a priori. I expected the 56\% Te shift to
propagate into a significant nF difference. It did not. That has
important implications for the surrogate, which I will discuss on the
next slide.

\subsection*{Figure Guidance}
If overlaid profiles (legacy vs.\ LXCat) for Te or ne are shown,
point to the region of maximum discrepancy. Then contrast with the
nF profiles, which should be nearly overlapping. The visual contrast
between ``big difference in Te'' and ``no difference in nF'' is the
key message.

\subsection*{Likely Questions}
\begin{itemize}
  \item \emph{Why does the discrepancy shrink for nF?}
    --- Fluorine density sits downstream of many competing reaction
    pathways; rate errors partially cancel.
  \item \emph{Should we trust LXCat more than legacy?}
    --- LXCat is more recent and community-vetted, but the solver
    validation (coming up) shows a mixed picture.
\end{itemize}

\subsection*{Transition}
The less-than-2\% nF change led me to do a formal dataset diagnosis.

\clearpage

%% ============================================================
\section{Slide 10 --- Dataset Diagnosis Surprise}
\label{sec:slide10}

\subsection*{Purpose}
Present the key finding: legacy and LXCat produce statistically
identical nF distributions. This simplifies the surrogate strategy.

\subsection*{What to Say}

I did a systematic comparison of the nF profiles from both
cross-section datasets across the full input space. The
standard-deviation ratio is 0.99$\times$ --- essentially unity.
Point by point, the nF profiles are indistinguishable within the
solver's own numerical noise.

This is genuinely surprising. The underlying plasma state is
different --- Te and ne shifted by 56\% and 69\% respectively ---
but the observable we care about is invariant. It means we do not
need to treat cross-section provenance as an input feature to the
surrogate. A single model trained on either dataset will generalise
to the other for nF predictions. It also means that any accuracy
differences between legacy-trained and LXCat-trained surrogates are
training artefacts, not physics.

I want to be careful here: this invariance is specific to nF. If we
needed to predict Te or ne, we would need separate models for each
cross-section dataset. But for our target observable, the two
datasets are interchangeable.

\subsection*{Figure Guidance}
If a parity plot (legacy nF vs.\ LXCat nF) is shown, point to the
tightness of the scatter around the 1:1 line. Mention the
0.99$\times$ standard-deviation ratio. If a residual histogram is
shown, point to its narrowness.

\subsection*{Likely Questions}
\begin{itemize}
  \item \emph{Could this change at different pressures or geometries?}
    --- Possible; we have not tested outside our current parameter
    range. Worth investigating.
  \item \emph{Does this mean the cross-section data does not matter?}
    --- No --- it matters for Te and ne. It just does not matter
    for nF, which is our prediction target.
\end{itemize}

\subsection*{Transition}
With this physics understood, I went back to the surrogate and ran
a proper architecture sweep.

\clearpage

%% ============================================================
\section{Slide 11 --- Architecture Sweep}
\label{sec:slide11}

\subsection*{Purpose}
Present the results of the architecture hyperparameter sweep.
Establish that E3 with separate heads is the winner.

\subsection*{What to Say}

I swept over several architecture variants: shared heads versus
separate heads, different trunk depths, hidden widths, and activation
functions. Each variant was trained multiple times with different
seeds, and I report the median test RMSE.

The winner is what I call E3 with separate heads. It achieves a test
RMSE of 0.00210 for a single model, which is already better than the
v4 architecture at 0.00292. The key insight is that separate output
heads let the network specialise its predictions for different spatial
regions --- the steep-gradient near-electrode zone versus the smoother
bulk. The shared-trunk design extracts common features from the inputs,
and then each head adapts those features to its own output region.

Wider and deeper networks did not help beyond a point. The
separate-heads design was the single biggest architectural win.

\subsection*{Figure Guidance}
If a bar chart or table of RMSE values by architecture is shown,
draw attention to the E3-separate-heads bar. Point to the gap between
it and the runner-up. If error bars are shown, note that the
improvement is statistically significant across seeds.

\subsection*{Likely Questions}
\begin{itemize}
  \item \emph{How many architectures were in the sweep?}
    --- Give the exact number.
  \item \emph{Did you use automated NAS?}
    --- No; the search space was small enough for a structured manual
    sweep.
  \item \emph{How much extra compute does separate heads cost?}
    --- Minimal at inference; the trunk is shared and the heads are
    small.
\end{itemize}

\subsection*{Transition}
Architecture is only half the story. The ablation tells us which
training ingredients actually matter.

\clearpage

%% ============================================================
\section{Slide 12 --- Ablation Study}
\label{sec:slide12}

\subsection*{Purpose}
Isolate the contribution of each training ingredient. Highlight the
surprising result that regularisation alone is useless.

\subsection*{What to Say}

I ablated two key ingredients: bias initialisation and regularisation
(weight decay). The full recipe combines both with learning-rate
warmup. Here are the numbers.

Bias initialisation alone gives an RMSE of 0.00314 --- that is 83\%
worse than the best ensemble, but it is the single most impactful
ingredient in isolation. It works because initialising the output
bias to the training-set mean gives the network a sensible starting
point; it only needs to learn the residual.

Regularisation alone gives 0.01888. That is essentially useless ---
it is worse than v3. Weight decay on its own constrains the network
too aggressively for our small dataset and compact architecture. It
prevents the model from fitting fine-grained spatial features.

The lesson for practitioners: not every ``best practice'' helps in
every context. Regularisation is only useful in combination with the
other techniques; in isolation, it is actively harmful. Bias
initialisation is the one thing you should always do.

\subsection*{Figure Guidance}
Point to the ablation table or bar chart. Highlight the
regularisation-only row and its surprisingly high RMSE --- this is
the most counterintuitive result. Contrast it with bias-only.
If the full-recipe result is shown, note that ensembling (next slide)
improves further.

\subsection*{Likely Questions}
\begin{itemize}
  \item \emph{What counts as ``full recipe''?}
    --- Bias init + weight decay + LR warmup, all together.
  \item \emph{Would more data change this picture?}
    --- Possibly; with a larger dataset, regularisation might become
    more beneficial.
  \item \emph{Did you try dropout instead of weight decay?}
    --- Yes/no; mention what was tried.
\end{itemize}

\subsection*{Transition}
With the best architecture and training recipe identified, here is
the production result.

\clearpage

%% ============================================================
\section{Slide 13 --- Production Ensemble Result}
\label{sec:slide13}

\subsection*{Purpose}
Present the headline number. This is the main result of the
update.

\subsection*{What to Say}

The production ensemble --- multiple E3-separate-heads models trained
with different random seeds --- achieves a test-set RMSE of 0.00154
in log\textsubscript{10} fluorine density. For SF\textsubscript{6}
density, the corresponding number is 0.00128.

To put 0.00154 in physical terms: $10^{0.00154} \approx 1.0035$,
so the typical prediction is within about 0.35\% of the true density
in linear space. This is 47\% better than the previous best single
model --- v4 at 0.00292. The improvement comes from the combination
of the E3 architecture, the full training recipe, and ensemble
averaging.

The ensemble also gives us uncertainty quantification for free. The
standard deviation across members flags inputs where the model is
less confident. In practice, the spread is tightest in the interior
of the training distribution and widens near the edges, which is
exactly the behaviour we want.

\subsection*{Figure Guidance}
If a parity plot (predicted vs.\ true log nF) is shown, point to the
tightness of the scatter around the 1:1 line. If a spatial profile
overlay is shown, point out that the predicted and true curves are
visually indistinguishable. Mention the nSF6 number if both species
are plotted.

\subsection*{Likely Questions}
\begin{itemize}
  \item \emph{How many ensemble members?}
    --- Give the exact count.
  \item \emph{Is 0.00154 the mean or median across test samples?}
    --- Clarify which statistic is reported.
  \item \emph{How does ensemble spread compare to RMSE?}
    --- Spread is typically comparable to or slightly larger than
    RMSE, so it is a useful, slightly conservative uncertainty
    estimate.
\end{itemize}

\subsection*{Transition}
Those are the ML metrics. But the surrogate is only as good as the
solver it learned from. Let me show the solver validation.

\clearpage

%% ============================================================
\section{Slide 14 --- Solver Validation}
\label{sec:slide14}

\subsection*{Purpose}
Establish credibility of the ground truth by comparing the solver
to published benchmarks for both cross-section datasets.

\subsection*{What to Say}

Before we trust the surrogate, we need to trust the solver. I
compared the solver's predictions against published literature for
two quantities: electron temperature and fluorine density.

For fluorine density, both datasets agree with literature to within
about 36\%. For electron temperature, the picture diverges: the
legacy cross-sections give about 5\% discrepancy relative to
literature, while LXCat shows 58\%. So LXCat, despite being more
recent and community-curated, does not uniformly improve agreement
with experiment.

This is an important nuance. It means the ``best'' cross-section
dataset depends on which observable you care about. For our
purposes --- predicting nF --- the solver is credible with either
dataset. The 36\% literature discrepancy sounds large, but the
literature values themselves carry significant experimental
uncertainty. What matters is that the solver is in the right
ballpark and internally consistent.

\subsection*{Figure Guidance}
If a comparison table is shown, point to the Te row and the nF row
separately. Emphasise that the conclusions are
observable-dependent --- this connects back to the dataset diagnosis
finding.

\subsection*{Likely Questions}
\begin{itemize}
  \item \emph{Which literature source?}
    --- Cite the specific papers used.
  \item \emph{Is 36\% error acceptable?}
    --- For this class of simulation, yes; it is within the range
    reported by other groups in the literature.
  \item \emph{Does the 58\% Te error under LXCat worry you?}
    --- It is worth investigating further, but since nF is invariant
    to the dataset choice, it does not affect our surrogate results.
\end{itemize}

\subsection*{Transition}
The solver is credible. Let me now show the spatial error structure
and mesh convergence.

\clearpage

%% ============================================================
\section{Slide 15 --- Spatial Error and Mesh Convergence}
\label{sec:slide15}

\subsection*{Purpose}
Show that (a) the solver is numerically converged and (b) the
surrogate's error is below the solver's own discretisation error.

\subsection*{What to Say}

I ran a mesh-convergence study by progressively refining the grid
and comparing nF profiles. The profiles converge to a fixed shape,
confirming we are in the asymptotic regime. The remaining nF
convergence error is 1.86\%, which is much smaller than the 36\%
physics uncertainty from the literature comparison. So the mesh is
not the bottleneck --- the physics uncertainty dominates.

More importantly, the surrogate's RMSE of 0.00154 --- corresponding
to 0.35\% linear error --- is well below the solver's own 1.86\%
mesh-convergence error. This means the surrogate is not adding
meaningful error on top of what the solver already introduces through
discretisation. We are faithfully reproducing the solver's output,
not a smoothed or aliased version of it.

\subsection*{Figure Guidance}
If a convergence plot (error vs.\ mesh nodes) is shown, point to the
plateau. If a spatial error map is shown, point to the
near-electrode regions where surrogate error is highest --- these
coincide with the steepest nF gradients, as expected.

\subsection*{Likely Questions}
\begin{itemize}
  \item \emph{What is the finest mesh you tested?}
    --- Give the node count.
  \item \emph{Is 1.86\% acceptable for the mesh error?}
    --- Yes; it is an order of magnitude below the physics uncertainty
    and an order of magnitude above the surrogate error. The error
    budget is dominated by cross-section uncertainty, not
    discretisation.
\end{itemize}

\subsection*{Transition}
Let me compare the neural network against simpler ML baselines.

\clearpage

%% ============================================================
\section{Slide 16 --- ML Baselines}
\label{sec:slide16}

\subsection*{Purpose}
Show that the neural network is not overkill by comparing against
Gaussian processes and polynomial regression.

\subsection*{What to Say}

I benchmarked two baselines. A Gaussian process regressor achieves an
RMSE of 0.006 --- roughly 4$\times$ worse than our ensemble at
0.00154. A polynomial regressor gives 0.039, which is essentially
unusable for our accuracy requirements.

The GP result is interesting. It is a respectable model that gives
calibrated uncertainty estimates and works well with small datasets.
But it scales poorly to larger training sets and does not handle the
high-dimensional spatial output as naturally as the neural network.
For our problem, the neural network is the right tool --- not because
simpler methods fail completely, but because they hit a performance
ceiling well above what we need.

The polynomial result is not surprising --- it confirms that the
input--output mapping has nonlinearities that a low-order polynomial
cannot capture. You need either a flexible kernel method (GP) or a
neural network.

\subsection*{Figure Guidance}
If a bar chart or table is shown, point to each baseline in turn.
Draw attention to the GP--NN gap (4$\times$) and the polynomial's
poor performance (25$\times$ worse than ensemble).

\subsection*{Likely Questions}
\begin{itemize}
  \item \emph{Did you tune the GP kernel?}
    --- Yes; Mat\'ern kernel with hyperparameter optimisation via
    marginal likelihood.
  \item \emph{What about XGBoost or random forests?}
    --- Tree methods handle multi-output poorly; we did not test
    them but would not expect them to beat the GP.
\end{itemize}

\subsection*{Transition}
Now let me quantify the speedup.

\clearpage

%% ============================================================
\section{Slide 17 --- Speedup}
\label{sec:slide17}

\subsection*{Purpose}
Quantify wall-clock speedup over both solver variants. This is one
of the two headline numbers (along with accuracy).

\subsection*{What to Say}

The surrogate achieves a 531$\times$ speedup over the legacy solver
and an 874$\times$ speedup over the LXCat solver. The LXCat solver
is slower because the updated cross-sections produce stiffer
chemistry that requires more sub-iterations in the implicit
time-stepper.

These speedups are measured end-to-end: solver time includes I/O
and post-processing; surrogate time includes input normalisation,
forward pass, and output de-normalisation. No cherry-picking. The
surrogate runs on a single CPU core with no GPU required --- for a
model this small, GPU launch overhead actually dominates, so CPU is
faster for single-sample inference.

In practical terms, the 874$\times$ number means that a parameter
sweep that would take a full workday with the LXCat solver finishes
in about 100 seconds with the surrogate. That is the difference
between ``run overnight'' and ``run while you wait.''

\subsection*{Figure Guidance}
If a bar chart of wall-clock times is shown, point to the three bars.
If a log scale is used, mention it. Emphasise that the surrogate bar
is barely visible on a linear scale --- that is the point.

\subsection*{Likely Questions}
\begin{itemize}
  \item \emph{Does this include batch inference?}
    --- Clarify whether the speedup is per-sample or amortised.
  \item \emph{What is the absolute inference time?}
    --- Milliseconds per profile on CPU.
  \item \emph{Would a GPU help for batch predictions?}
    --- Yes, for large batches; for single-sample, CPU wins.
\end{itemize}

\subsection*{Transition}
Let me step back and rank all the things that mattered.

\clearpage

%% ============================================================
\section{Slide 18 --- Computational Cost}
\label{sec:slide18}

\subsection*{Purpose}
Show that the experimental campaign was thorough and that the
compute budget was spent deliberately.  Demonstrate awareness of
cost---important for credibility and reproducibility.

\subsection*{What to Say}

Before I rank the findings, let me give you the full accounting of
compute time.  The entire experimental campaign consumed 97 hours
of training across 13 experiment categories.  That breaks down into
three tiers.

The core experiments---architecture sweep, ablation study, and
production ensemble---took about 41 hours and produced all of the
positive results I have shown you.  The supplementary
experiments---transfer learning, mixed-physics training, and the
Te auxiliary head---consumed 54 hours, more than the core work, and
all three yielded negative results.  The remaining 3 hours cover
fast validation tasks like ML baselines, mesh convergence, and
literature comparison.

I want to draw your attention to the fact that the supplementary
experiments dominate the compute budget.  The ablation study alone
took 22 hours because five configurations were each trained with
three seeds.  Transfer learning and mixed-physics training each
required multiple full training runs.  None of these experiments
improved the production surrogate, but they preemptively answer
the obvious reviewer questions: ``Did you try transfer learning?''
Yes---no benefit.  ``Can one model handle both cross-section sets?''
No---it hurts both.  ``Does predicting Te help?'' No---it hurts nF
by 11\%.

All training ran on Apple MPS for the GPU experiments and on CPU for
ablation and supplementary runs.  Multiple experiments ran in
parallel to reduce wall-clock time.

\subsection*{Figure Guidance}
The slide shows a compact table grouped into Core, Supplementary,
and Fast Validation tiers.  Point to the hours column for the
supplementary tier and emphasise that 55\% of compute went to
negative results.

\subsection*{Likely Questions}
\begin{itemize}
  \item \emph{Why did the supplementary experiments take so long?}
    --- Each required full training to convergence with multiple
    seeds; the LXCat dataset training is inherently expensive.
  \item \emph{Could you have skipped the negative experiments?}
    --- In hindsight, yes, but reviewers will ask about transfer
    learning and multi-task learning.  Having the results ready
    strengthens the paper.
  \item \emph{What hardware was used?}
    --- Apple MPS (GPU) for architecture sweep and ensemble; CPU for
    ablation and supplementary experiments.
\end{itemize}

\subsection*{Transition}
Now let me rank everything that mattered.

\clearpage

%% ============================================================
\section{Slide 19 --- What Mattered Most}
\label{sec:slide19}

\subsection*{Purpose}
Synthesise findings into a ranked list. Include both positive and
negative results.

\subsection*{What to Say}

If I had to rank the factors that mattered most for getting from a
mediocre surrogate to a production-quality one, here is my list.

Number one: working in log space. This makes the regression problem
well-conditioned. Without it, the network struggles with the dynamic
range. Number two: the separate-heads architecture --- the biggest
architectural win. Number three: bias initialisation --- the most
impactful single training trick, as confirmed by the ablation.
Number four: ensembling, which smooths out seed variance and provides
uncertainty estimates.

Number five is just as important: what did \emph{not} work.
Regularisation alone was useless --- RMSE of 0.01888. Transfer
learning from legacy to LXCat showed no benefit: 0.00375 versus
0.00385 from scratch. Mixed-physics training, combining both
datasets, hurt both tasks: 0.006 for legacy, 0.023 for LXCat.
Predicting Te as an auxiliary output did not help nF either --- 0.004
in both cases, with added complexity. PINNs were a complete dead end.

I think the negative results are just as valuable as the positive
ones. They save future researchers from going down dead ends that
look promising on paper.

\subsection*{Figure Guidance}
If a ranked list or waterfall chart is shown, walk through it top to
bottom. Dwell on the negative results at the bottom --- the
supervisor will find these most informative.

\subsection*{Likely Questions}
\begin{itemize}
  \item \emph{Would more data change the ranking?}
    --- Possibly; with a larger dataset, regularisation might become
    more important.
  \item \emph{Is the ranking specific to this problem?}
    --- Log-space is likely general; the rest may be
    problem-dependent.
\end{itemize}

\subsection*{Transition}
Let me be honest about the limitations.

\clearpage

%% ============================================================
\section{Slide 20 --- Limitations}
\label{sec:slide20}

\subsection*{Purpose}
Acknowledge what the current work cannot do. Demonstrate
intellectual honesty and identify where effort should go next.

\subsection*{What to Say}

I want to be upfront about what this surrogate cannot do yet. First,
it predicts steady-state profiles only --- no transient dynamics.
If you need the time evolution of the discharge, you still need the
solver. Second, the model is trained on a fixed parameter range; I
have not tested extrapolation beyond the training domain, and I would
not trust it there without additional validation.

Third, the mesh convergence error is 1.86\%, which means the solver
itself introduces that much uncertainty before the surrogate even
enters the picture. The surrogate cannot be more accurate than its
training data. Fourth, the LXCat Te discrepancy of 58\% versus
literature is concerning --- if Te accuracy ever becomes important,
we will need to revisit the cross-section data or the solver's
treatment of electron energy balance.

Finally, the negative results on transfer learning and mixed-physics
training suggest that combining data from different cross-section
datasets is nontrivial. If new datasets emerge in the future, we
cannot simply throw them into the existing training set --- we will
need separate models or a more sophisticated conditioning approach.

\subsection*{Figure Guidance}
No complex figure expected. If a bullet-point summary of limitations
is shown, walk through each point and give the one-line justification
for why it matters.

\subsection*{Likely Questions}
\begin{itemize}
  \item \emph{How would you handle extrapolation?}
    --- Use the ensemble spread as a warning flag; retrain with
    additional solver runs in the new regime.
  \item \emph{Is transient prediction feasible with this approach?}
    --- Possibly with a sequence model (RNN, transformer), but that
    is future work.
\end{itemize}

\subsection*{Transition}
Given those limitations, here is what I plan to do next.

\clearpage

%% ============================================================
\section{Slide 21 --- Next Steps}
\label{sec:slide21}

\subsection*{Purpose}
Summarise the key takeaways and outline the near-term plan. End
on a concrete, actionable note.

\subsection*{What to Say}

Here is what I plan to work on next.  First, I want to extend the
surrogate to time-dependent profiles, probably using a recurrent or
transformer-based decoder.  Second, I want to investigate active
learning to minimise the number of solver runs needed for
training --- right now we do a brute-force grid, and there may be
smarter sampling strategies.  Third, I want to deploy the surrogate
in an interactive tool where an engineer can adjust operating
conditions and see the fluorine profile update in real time.

Phase 2 on the physics side would be integrating BOLSIG+ for the
EEDF calculation, then a PINN-based Boltzmann solver, and eventually
PIC-MCC validation.  These are longer-term goals, but the surrogate
infrastructure built here carries over directly.

\subsection*{Figure Guidance}
If a roadmap or bullet list is shown, walk through each item briefly.
Point to the benchmark plot on the right as evidence that the
surrogate infrastructure is mature enough to build on.

\subsection*{Likely Questions}
\begin{itemize}
  \item \emph{What is the timeline for transient modelling?}
    --- Exploratory; depends on data availability.
  \item \emph{How would active learning work in practice?}
    --- Use ensemble uncertainty to identify high-uncertainty regions,
    run solver at those conditions, and retrain.
\end{itemize}

\subsection*{Transition}
Let me close with the summary.

\clearpage

%% ============================================================
\section{Slide 22 --- Conclusion}
\label{sec:slide22}

\subsection*{Purpose}
Deliver the three-point summary and the key insight.  End on a
memorable, concrete note.

\subsection*{What to Say}

Let me wrap up with three takeaways.

First, the supervised ensemble surrogate works: 0.00154 RMSE in log
nF, 0.35\% linear error, 531--874$\times$ speedup.  This is
production-ready for the use cases we discussed.

Second, the cross-section dataset choice does not affect nF
predictions --- the 0.99$\times$ standard-deviation ratio was an
unexpected but useful finding.  LXCat and legacy produce
statistically identical nF distributions despite 56\% shifts in
electron temperature.

Third, several popular ML techniques --- PINNs, transfer learning,
multi-task learning --- did not help here, and documenting those
failures is important.  The full campaign consumed 97 hours of
compute, with 55\% going to supplementary experiments that all
yielded negative results.

The single most important lesson: bias initialisation accounts for
83\% of the improvement.  Training recipe dominates architecture,
which dominates data source.  Recipe is everything.

That is everything I have.  Thank you --- happy to take questions.

\subsection*{Figure Guidance}
The slide shows a numbered summary and a boxed key insight.  Do not
re-read numbers from the screen; the supervisor has them visible.
Let the boxed statement land on its own.

\subsection*{Likely Questions}
General Q\&A.  Be prepared for:
\begin{itemize}
  \item Deep dives into the PINN failure mode.
  \item Questions about generalisation to other gas mixtures.
  \item Requests for the code or dataset.
  \item How the uncertainty estimate would be used in practice.
  \item Timeline for the next-steps items.
\end{itemize}

\subsection*{Transition}
If questions arise about transfer learning, Te auxiliary outputs,
mixed-physics training, or PINN details, advance to the backup
slides.

\clearpage

%% ============================================================
\section*{Backup Slides}

%% ============================================================
\subsection*{Backup A --- PINN Details}

\textbf{Purpose:} Provide additional detail on the PINN experiments
if the supervisor wants to dig deeper.

\textbf{What to Say:}

The data-only PINN used the same network architecture as our
supervised baseline but trained through the PINN framework with
automatic differentiation and collocation points at the mesh nodes.
It converged to the same loss level as standard training, but
required 23,000$\times$ more wall-clock time. The overhead comes
from evaluating the PDE residual at every collocation point at every
training step, even though we are not using it in the loss --- it
is purely framework overhead.

The PDE-constrained variant added the drift-diffusion equations
as a soft penalty. I tried multiple loss-weighting strategies:
fixed weights from 0.01 to 100, self-adaptive weights, and
NTK-based balancing. None stabilised training. The fundamental
issue is that the PDE residual involves stiff source terms with
rate constants spanning 10+ orders of magnitude, which creates
conflicting gradient directions that the optimiser cannot reconcile.

\textbf{Likely Questions:}
\begin{itemize}
  \item \emph{Would causal training (Krishnapriyan et al.) help?}
    --- Possibly for the temporal dimension, but we are predicting
    steady state, so causality is not the issue.
  \item \emph{What about domain decomposition for PINNs?}
    --- An interesting idea but adds significant implementation
    complexity for uncertain benefit.
\end{itemize}

\bigskip

%% ============================================================
\subsection*{Backup B --- Transfer Learning}

\textbf{Purpose:} Show that pre-training on legacy data and
fine-tuning on LXCat data provides no benefit.

\textbf{What to Say:}

I hypothesised that the legacy dataset could serve as a pre-training
source for the LXCat surrogate. I pre-trained on legacy data, froze
the trunk, and fine-tuned the heads on LXCat data. The result:
0.00375 RMSE, compared to 0.00385 for training on LXCat from scratch.
The difference is within noise.

This is entirely consistent with the dataset diagnosis. If the nF
distributions are statistically identical, there is nothing meaningful
to transfer. The 0.99$\times$ standard-deviation ratio tells us the
two datasets carry the same information about nF, so pre-training on
one does not give the network a head start on the other. I also tried
unfreezing the trunk during fine-tuning --- results were similar.

\textbf{Likely Questions:}
\begin{itemize}
  \item \emph{What if the datasets were more different?}
    --- Then transfer might help, but the diagnosis tells us they
    are not different enough.
  \item \emph{Could transfer help with a completely new gas mixture?}
    --- Possibly; that would be a genuinely different distribution.
\end{itemize}

\bigskip

%% ============================================================
\subsection*{Backup C --- Te Auxiliary Output}

\textbf{Purpose:} Show that multi-task learning with Te does not
improve nF accuracy.

\textbf{What to Say:}

Multi-task learning is a standard regularisation technique: if
electron temperature shares the same input features as fluorine
density, predicting both should improve the shared representation.
In practice, the two-output model achieves an nF RMSE of 0.004 ---
identical to the single-output model. The Te head adds parameters and
training complexity with no benefit to the quantity we care about.

My interpretation is that nF and Te have sufficiently different
functional forms that sharing a trunk forces compromises. The
separate-heads architecture partially mitigates this, but the net
effect is still neutral at best. The added training complexity is
not worth it.

\textbf{Likely Questions:}
\begin{itemize}
  \item \emph{What loss weighting did you use?}
    --- Swept over several ratios; none improved nF.
  \item \emph{Could Te help as an input feature?}
    --- Interesting idea, but Te is a solver output, not known a
    priori.
\end{itemize}

\bigskip

%% ============================================================
\subsection*{Backup D --- Mixed-Physics Training}

\textbf{Purpose:} Show that combining legacy and LXCat data into
one training set hurts both tasks.

\textbf{What to Say:}

I tried training a single surrogate on the combined legacy + LXCat
dataset, adding a binary indicator feature for dataset provenance.
The result was worse for both subsets: 0.006 RMSE on legacy data
(versus 0.00154 for the specialist model) and 0.023 on LXCat data.
The network cannot simultaneously serve two masters.

This is a case where the 0.99$\times$ nF similarity is misleading.
The distributions are similar but not identical, and the small
differences act as label noise when mixed. The network wastes
capacity reconciling conflicting labels at nearly identical input
points. The practical recommendation: train separate surrogates for
each cross-section dataset. The computational cost is minimal, and
the accuracy benefit is large.

\textbf{Likely Questions:}
\begin{itemize}
  \item \emph{Did you try FiLM layers or other conditioning?}
    --- No; worth exploring if cross-dataset generalisation were
    a hard requirement.
  \item \emph{What about a mixture-of-experts approach?}
    --- That would effectively learn separate models, which is
    what we already recommend.
\end{itemize}

\clearpage

%% ============================================================
\section*{Q\&A Preparation}

This section collects the most likely questions across all slides,
organised by topic. Review this before the meeting.

\subsection*{Physics and Solver}

\begin{enumerate}[1.]
  \item \emph{Is the solver truly 2D or 1D?} ---
    The solver is 2D (TEL geometry). The surrogate predicts the
    spatial profile along the gap axis, which is the
    safety-relevant dimension.

  \item \emph{Why does the LXCat Te discrepancy not propagate to nF?} ---
    Fluorine density is downstream of many competing reactions whose
    rate errors partially cancel. Te affects individual rates, but the
    integrated effect on nF is buffered.

  \item \emph{What is the 36\% literature error on F density?} ---
    Both legacy and LXCat solvers agree with published experimental /
    computational benchmarks to within $\sim$36\%. The literature values
    themselves carry experimental uncertainty.

  \item \emph{Is 1.86\% mesh convergence error acceptable?} ---
    Yes. It is one order of magnitude below the physics uncertainty
    (36\%) and one order of magnitude above the surrogate error
    (0.35\%). The error budget is dominated by cross-section
    uncertainty, not discretisation.
\end{enumerate}

\subsection*{Machine Learning Methodology}

\begin{enumerate}[1.]
  \item \emph{Why not use PINNs?} ---
    Data-only PINNs are 23,000$\times$ slower than supervised
    training. PDE-constrained PINNs diverge due to stiff chemistry.
    With abundant solver data, supervised learning is faster, simpler,
    and more accurate.

  \item \emph{Why does regularisation alone hurt?} ---
    With a small dataset and compact network, weight decay prevents the
    model from fitting fine-grained spatial features. Bias init is far
    more impactful (0.00314 vs.\ 0.01888).

  \item \emph{How many ensemble members?} ---
    Give the exact count used in the production model.

  \item \emph{Why separate heads instead of a shared output layer?} ---
    The nF profile has steep gradients near the electrodes and smooth
    behaviour in the bulk. Separate heads let the network specialise
    for each regime. The architecture sweep confirms a win:
    0.00210 (E3 separate) vs.\ 0.00292 (v4 shared).

  \item \emph{Did you try transformers, DeepONet, or FNO?} ---
    Not yet. The problem is small enough that an MLP with separate
    heads works well. More complex architectures are possible future
    directions.

  \item \emph{How sensitive is the result to hyperparameters?} ---
    The v1$\to$v4 progression shows that training recipe matters more
    than fine-tuning. The ablation isolates which ingredients drive the
    improvement.
\end{enumerate}

\subsection*{Negative Results}

\begin{enumerate}[1.]
  \item \emph{Why did transfer learning not help?} ---
    The dataset diagnosis shows legacy and LXCat nF distributions have
    a 0.99$\times$ std ratio --- they carry the same information, so
    there is nothing to transfer. Result: 0.00375 vs.\ 0.00385.

  \item \emph{Why did mixed-physics training hurt?} ---
    Small differences between datasets act as label noise. The network
    wastes capacity reconciling conflicting labels: 0.006 legacy,
    0.023 LXCat, both worse than separate models.

  \item \emph{Why did the Te auxiliary output not help nF?} ---
    nF and Te have different functional forms; sharing a trunk forces
    compromises. Both single- and two-output models give nF RMSE of
    0.004.

  \item \emph{Are these negative results problem-specific?} ---
    Partly. The PINN failure is likely general for stiff chemistry.
    The transfer/mixed-physics results are specific to the near-identical
    dataset pair. The Te auxiliary result may generalise to other
    multi-output problems where outputs have very different functional
    forms.
\end{enumerate}

\subsection*{Practical Deployment}

\begin{enumerate}[1.]
  \item \emph{Can I use the surrogate for conditions outside the
    training range?} ---
    Not recommended without validation. The ensemble spread provides a
    warning flag for out-of-distribution inputs. Retrain with
    additional solver runs in the new regime.

  \item \emph{How would the uncertainty estimate be used in practice?} ---
    Flag predictions where ensemble spread exceeds a threshold.
    Trigger a solver run for those cases to validate or augment the
    training set.

  \item \emph{What is the deployment plan?} ---
    An interactive tool where an engineer adjusts operating conditions
    (pressure, voltage, gap) via sliders and sees the nF profile update
    in real time. Millisecond inference on CPU makes this feasible.

  \item \emph{Can I get the code or the trained models?} ---
    Discuss availability and any IP considerations.
\end{enumerate}

\clearpage

%% ============================================================
\section*{Quick-Reference Number Sheet}

For convenience during Q\&A, all key numbers in one place:

\begin{center}
\begin{tabular}{l l}
\hline
\textbf{Metric} & \textbf{Value} \\
\hline
Ensemble nF RMSE (log\textsubscript{10}) & 0.00154 ($\sim$0.35\% linear) \\
Ensemble nSF6 RMSE (log\textsubscript{10}) & 0.00128 \\
Legacy v4 single-model RMSE & 0.00292 \\
Improvement over v4 & 47\% \\
\hline
\multicolumn{2}{l}{\emph{Surrogate progression}} \\
v1 & 0.133 \\
v2 & 0.111 \\
v3 & 0.023 \\
v4 & 0.003 \\
\hline
Architecture sweep winner & E3 separate heads (0.00210) \\
\hline
\multicolumn{2}{l}{\emph{Ablation}} \\
Bias only & 0.00314 (+83\% from best) \\
Reg only & 0.01888 (useless) \\
\hline
Speedup (legacy solver) & 531$\times$ \\
Speedup (LXCat solver) & 874$\times$ \\
\hline
\multicolumn{2}{l}{\emph{Literature comparison}} \\
F density error & 36\% \\
Te error (legacy) & 5\% \\
Te error (LXCat) & 58\% \\
\hline
\multicolumn{2}{l}{\emph{LXCat physics impact}} \\
Te change & +56\% \\
ne change & $-$69\% \\
nF change & $<$2\% \\
\hline
\multicolumn{2}{l}{\emph{ML baselines}} \\
GP RMSE & 0.006 \\
Polynomial RMSE & 0.039 \\
\hline
\multicolumn{2}{l}{\emph{PINN}} \\
Data-only overhead & 23,000$\times$ \\
PDE-constrained & diverges \\
\hline
\multicolumn{2}{l}{\emph{Negative results}} \\
Transfer learning & 0.00375 vs.\ 0.00385 (no benefit) \\
Mixed-physics legacy & 0.006 \\
Mixed-physics LXCat & 0.023 \\
Te auxiliary nF RMSE & 0.004 vs.\ 0.004 \\
\hline
Mesh convergence error (nF) & 1.86\% \\
\hline
\end{tabular}
\end{center}

\end{document}
